% Transcription — fonds Alexandre Grothendieck, shelfmark 58, batch 2 (pages 21-40).
% Established from the facsimile archives/batches/58/batch-02.pdf (PDF page k =
% archivists' page 20+k, no cover sheet), cut from a copy of 58.pdf fetched
% through the project's own relay
% (https://grothendieck-relay.onrender.com/source/58.pdf); tiles in
% archives/tiles/58/p21-p40, re-cropped at 300-400 dpi with pdftoppm. Per the
% skill transcribe-grothendieck. The folder is sixty-six pages in four batches
% (1-20, 21-40, 41-60, 61-66), transcribed in parallel; this is the second,
% written after batch 1 (whose notation it keeps) and without sight of 3-4.
%
% Pass: Opus 5.5 (claude-opus-5-5), 2026-09-26 — first pass, unchecked against the pages by a human.
%
% Pages skipped: 21, 23, 25 — typed leaves of a Bourbaki draft (« Algèbre
% commutative, Rédaction n° 355, Chapitre IV — Décomposition primaire »),
% the sheets whose backs carry pp. 22, 24, 26, with nothing in his hand;
% skipped by the settled rule for typed versos. 28 — a blank cream sheet.
% Page 27 is a grey cover inscribed top right « [struck word] couple de
% Lefschetz-Grauert (X, Y) »; it takes no \page marker and its words become
% the section heading, with a \note.
%
% Pages summarised: none. Pages transcribed: 22, 24, 26 (blue ink on the
% typed versos) and 29-40 (pencil and ink on thin grey paper). Pages 29-30
% (an « Introduction ») are a fast palimpsest with long cancelled passages
% and marginal additions; the formulas and a few phrases carry, most of the
% prose is left \ill{}. Pages 31-40 are written much more fairly and read
% largely end to end.
%
% His pagination: none visible. His numbered runs: pp. 22-26 « Théorème 2 »,
% « Lemme 5 », Cor. 1, 2, 3 « au th. 2 », referring back to Lemme 1 and
% Lemme 4 (batch 1 p. 20 has a Lemme 4). Page 29 is headed « III » (perhaps
% a chapter number, the next sign illegible) and « Introduction »; it cites
% « § 4 et 5 », « N° 1 », « N° 2 ». Page 31 is a plan in nine numbered
% items (3-4 cancelled and renumbered). Pages 32-40 follow items 5-7 of that
% plan: « Fermés de Lefschetz-Grauert », sections 1) Définitions, remarques
% (p. 32), 2) Premières propriétés (p. 37), 3) Théorème de permanence
% (p. 40); Remarques 1, 2, 3 (pp. 33, 33, 35), unnumbered Lemmes,
% Propositions and Corollaires, Théorème 1 and Théorème 2 (p. 40, the latter
% breaking off at the foot of the page).
%
% What the batch is about. Pp. 22-26: Théorème 2 — f : X -> Y proper, Y
% integral and (geometrically) unibranch at y_0, components of X dominating
% Y, generic fibre X_1 with components of dim >= k+1 and connected in dim
% >= k; then the special fibre X_0 is connected in dim >= k. Reduced via
% Lemme 4 (batch 1) to Lemme 5 on the punctured local spectra of X_0, proved
% by dim O_{X,x} = dim O_{Y,y_0} + dim X_1 and the local theorem of batch 1
% applied after dividing by n equations; corollaries: geometric variants, f
% proper and open, and X in P^n_k with successive hyperplane sections.
% Pp. 29-31: introduction and plan of a chapter comparing R^i f and
% R^i \hat f (formal functions, existence), towards Lefschetz-type theorems
% « via the formal completion X_{/Y} » (Grauert 1958, Serre duality), and the
% plan: finiteness, local finiteness theorem, the two special cases (local
% rings, projective schemes), Lefschetz-Grauert closed sets, pi_0, pi_1, Pic,
% permanence, purity and Lefschetz for pi_1, quasi-factorial rings and
% Lefschetz for Pic. Pp. 32-40: definition of (LG)_n for a closed subset Y of
% a locally noetherian X (H^i(X,F) -> H^i(X_{/Y},F_{/Y}) bijective for i <= n,
% F locally free), LG (faible), LG stricte / effective (lim over U of L(U) ->
% L(X_{/Y}) an equivalence); Remarque 2 reduces (LG)_n to a generating family
% (L_alpha) by induction and the five lemma; Remarque 3 and a Lemme on Ass;
% boxed « Conclusion pratique » for X quasi-projective. Then pi_0 (surjective/
% bijective), pi_1 (surjective; iso with lim pi_1(U) under LG strict), étale
% covers (fully faithful restriction; extension to a neighbourhood), Pic
% (injective/iso), a Corollaire for Y a Cartier divisor with Lambda = L(Y)
% and vanishing of H^1, H^2(Y, Lambda^{-n}), its Lemme on invertible sheaves
% on \hat X, and the start of « Théorème de permanence » (Théorème 1: (LG)_n
% passes to X' projective and flat over X quasi-projective).
%
% Notation fixed once for the batch: X_{/Y} for his formal completion (he
% writes X with a slanted /Y subscript), F_{/Y} likewise; \hat{F}, \hat{X}
% where he puts a hat; (\mathrm{LG})_n, (\mathrm{LG}) as he writes them in
% bold capitals; \mathcal{L}(U) for his script L (category of locally free
% Modules); \Lambda for his capital lambda; \mathcal{J} for the ideal of Y;
% \underline{L}(Y) as he underlines it; \mathcal{O}_X underlined in his hand
% is set plain. Batch 1's X', Y', \mathfrak{r}A are kept where they recur.
%
% Readings to check first: the whole of pp. 29-30; the opening of p. 22
% (« le point générique », « unibranche », the struck words); the formula
% dim >= K+1 = (n+k)+1 on p. 26 and its capital K; the items 3-4 of the plan
% (p. 31, « Théorie des ... locaux », the circled insertion); « faible » and
% « stricte » on p. 32 and « effective » on p. 35; the third member (LG)_X
% of the boxed conclusion p. 36; the struck H^1 in (ii) of the Corollaire
% p. 39; the cancelled first lines of p. 40.

\documentclass[11pt,a4paper]{article}
% Le préambule commun à toutes les transcriptions du fonds Grothendieck.
%
% Il tient en une idée : **l'incertitude se compose, elle ne se résout pas.**
% Une transcription qui lisse un mot douteux a détruit la seule chose pour
% laquelle on la faisait — un lecteur doit pouvoir savoir, à chaque endroit,
% ce qui était sur le feuillet et ce qui a été ajouté. Les six macros qui
% suivent sont donc l'appareil critique, et non de la décoration.
%
% Les mêmes commandes sont reconnues par `scripts/render.mjs`, qui produit la
% vue de lecture du site. Une macro ajoutée ici doit l'être aussi là-bas,
% faute de quoi le rendu échouera bruyamment — ce qui est voulu : un rendu
% silencieusement faux est pire qu'un rendu absent.

\ProvidesPackage{grothendieck}[2026/08/08 Transcriptions du fonds Grothendieck]

\usepackage{amsmath,amssymb,amsthm}
\usepackage{mathrsfs}
\usepackage{stmaryrd}
% Les diagrammes commutatifs sont partout dans ce fonds : c'est la langue
% même de la géométrie algébrique de Grothendieck, et une transcription qui
% les rendrait en prose ne transcrirait rien.
\usepackage{tikz-cd}
% Français par défaut — c'est la langue des feuillets — mais l'anglais est
% chargé aussi : la traduction et le résumé sont des documents anglais, et la
% typographie française (espaces avant les deux-points) y serait une faute.
% Ces documents basculent par \selectlanguage{english} en tête de corps.
\usepackage[english,french]{babel}
\usepackage{fontspec}
\usepackage{geometry}
\geometry{a4paper,margin=25mm,marginparwidth=28mm}
\usepackage{marginnote}
\usepackage{xcolor}
% ulem, not soul. `soul`'s \st and \ul are built on a character-by-character
% scan that XeLaTeX's font machinery breaks: the document fails with an
% undefined control sequence rather than degrading. ulem does the same two jobs
% and is Unicode-engine safe. `normalem` stops it hijacking \emph, which the
% transcriptions use for Grothendieck's own emphasis.
\usepackage[normalem]{ulem}
\usepackage{hyperref}
\hypersetup{colorlinks=true,linkcolor=black,urlcolor=black,pdfborder={0 0 0}}

% --- Métadonnées du lot -----------------------------------------------------
% Reprises telles quelles par le script de rendu, qui les affiche en tête de
% la vue de lecture. Elles fixent ce que le fichier prétend couvrir : sans
% elles, une transcription flotte sans point d'attache dans un fonds de seize
% mille feuillets.

\newcommand{\folder}[1]{\gdef\grothFolder{#1}}
\newcommand{\batch}[1]{\gdef\grothBatch{#1}}
\newcommand{\leaves}[2]{\gdef\grothFirst{#1}\gdef\grothLast{#2}}
\newcommand{\foldertitle}[1]{\gdef\grothFoldertitle{#1}}
\gdef\grothFolder{?}\gdef\grothBatch{?}\gdef\grothFirst{?}\gdef\grothLast{?}\gdef\grothFoldertitle{}

% --- Le résumé --------------------------------------------------------------
%
% Il ouvre la lecture modernisée, sous le titre « Résumé », et il est composé
% en petit. Ce n'est pas un ornement : c'est le seul passage du document qui
% ne soit pas des mathématiques, et un lecteur doit pouvoir le reconnaître
% avant de l'avoir lu — puis le sauter s'il connaît déjà le sujet, ou s'y
% arrêter s'il ne le connaît pas. Le filet qui le suit marque la reprise du
% corps.

\newenvironment{resume}
  {\small\setlength{\parskip}{0.45em}}
  {\par\medskip\hrule\medskip}

% --- Mots-clés ---------------------------------------------------------------
%
% En anglais, en clôture du résumé de la lecture modernisée : le vocabulaire
% moderne sous lequel chercher ce que les feuillets construisent. Le manifeste
% du site (`npm run manifest`) les extrait pour étiqueter la cote entière —
% cette ligne est la source unique des tags, il n'y a pas de fichier à côté.
\newcommand{\keywords}[1]{\par\smallskip\noindent
  {\footnotesize\textsc{Keywords} --- \textit{#1}}\par}

% --- L'appareil critique ----------------------------------------------------

% Le numéro de feuillet, en marge. C'est l'ancre qui permet de régler un
% désaccord sur une formule en regardant *un* feuillet plutôt qu'une chemise
% de six cents. Le site s'en sert aussi pour tourner les pages du fac-similé
% quand on fait défiler la transcription.
\newcommand{\leaf}[1]{%
  \marginnote{\footnotesize\color{gray!70}\textsf{#1}}%
  \label{leaf:#1}\ignorespaces}

% Illisible : on ne devine pas. Un mot inventé qui se lit comme les autres est
% le pire résultat possible.
\newcommand{\ill}{\textcolor{red!60!black}{[\,\ldots\,]}}

% Lecture douteuse : le mot est donné, et signalé comme incertain.
\newcommand{\uncertain}[1]{\uline{#1}}

% Ajout de l'éditeur — une lettre manquante, un mot sous-entendu.
\newcommand{\add}[1]{\textcolor{blue!50!black}{[#1]}}

% Biffé par Grothendieck lui-même. Ce qu'il a rayé fait partie du document :
% une rature dit souvent où la pensée a changé de direction.
\newcommand{\struck}[1]{\sout{#1}}

% Note du transcripteur, sur l'état matériel du feuillet.
\newcommand{\note}[1]{\par\smallskip{\small\color{gray!60!black}\textit{[#1]}}\par\smallskip}

% Note marginale de Grothendieck, distincte du corps du texte.
\newcommand{\marginal}[1]{\par\smallskip\noindent\hspace{1em}%
  {\small\color{gray!60!black}#1}\par\smallskip}

% --- Titre ------------------------------------------------------------------

\AtBeginDocument{%
  \begin{center}
    {\large\bfseries Fonds Alexandre Grothendieck --- cote n\textsuperscript{o}~\grothFolder}\\[2pt]
    {\itshape\grothFoldertitle}\\[4pt]
    {\small feuillets \grothFirst--\grothLast\ (lot \grothBatch)}\\[2pt]
    {\footnotesize\color{gray!60!black}Transcription établie d'après le fac-similé numérisé
     par l'Université de Montpellier.\\
     Les lectures incertaines sont soulignées, les illisibles notées [\,\ldots\,],
     les ajouts entre crochets.}
  \end{center}
  \vspace{1em}}

\endinput


\folder{58}
\batch{2}
\pages{21}{40}
\dating{[à partir de 1964]}
\watermark{Édition de démonstration}
\foldertitle{Théorie de " Lefschetz-Grauert ". Applications en π₁, à Pic etc… : notes manuscrites (s.d.)}

\begin{document}

\page{22}
\emph{Théorème 2.} Soit $f\colon X \to Y$ un morphisme \struck{\ill{}} propre de préschémas \uncertain{localement} noethériens, avec $Y$ intègre, $y_0 \in Y$, $y_1$ \struck{\ill{}} le point générique de $Y$, \uncertain{$X_0$, $X_1$ les fibres}. Supposons :
\note{Pages 22, 24 et 26 : encre bleue, au dos de feuillets dactylographiés ; l'énoncé est encadré à gauche d'un trait vertical.}

(a) $Y$ \uncertain{est} [\uncertain{géométriquement}] unibranche en $y_0$,
\marginal{\struck{\ill{}} \ill{}}
et les composantes irréductibles de $X$ dominent $Y$.
\note{« (a) » est récrit sur un premier signe ; un ajout encadré et biffé, en haut à droite, est illisible.}

(b) les composantes irréductibles de $X_1$ sont de dim $\geq k$, \struck{(b) \ill{} \ill{} \ill{}} et $X_1$ \uncertain{est} connexe en dim $\geq k$.
\note{Le « $k$ » de la première condition de (b) pourrait être « $k+1$ » ; le tracé est serré en bout de ligne.}

Alors (les composantes irréductibles de $X_0$ sont de dim $\geq k+1$) et $X_0$ est connexe en dim $\geq k$.

L'assertion entre parenthèses est bien connue par la théorie de la dimension, et n'utilise pas l'hypothèse unibranche. Celle-ci sert uniquement à appliquer le \emph{th. de connexion}, qui assure que $X_0$ est connexe.

\page{24}
Le Lemme~4 nous \uncertain{ramène} alors à vérifier le

\emph{Lemme 5} (\uncertain{variante} \ill{} \ill{} th.). Sous les conditions du th., mais \struck{\ill{}} $X_0$ \struck{\ill{}} \uncertain{unibranche}, \uncertain{pour} $x \in X_0$, \struck{\ill{}} \ill{} $x$ de $X_0$, $\operatorname{Spec}(\mathcal{O}_{X_0,x})'$ est connexe en dim $\geq k$.
\marginal{\uncertain{+ théorème de dim.}}
\marginal{dim $\overline{x}$ \ill{} tel que dim $\overline{x} \geq k$, \uncertain{que} $\operatorname{Spec}(\mathcal{O}_{X_0,x})'$ \ill{}}
\note{L'ajout « + théorème de dim » est encerclé à droite ; la note marginale gauche, oblique, est reliée à l'énoncé.}

En effet, le Lemme~1 \ill{} \ill{} \uncertain{nous ramène} facilement au cas où \ill{} l'hypothèse sur $X_1$, \ill{} \ill{} :
$X_1$ est irréductible \struck{de dim $\geq k+1$}, de dim $\geq k+1$,
$X$ irréductible, \ill{} \ill{} de $Y$, donc $=$ \uncertain{intègre} \uncertain{comme} $Y$.
Les \ill{} \ill{} \ill{} \ill{} \struck{\ill{}} \ill{}
\struck{\ill{} \ill{}} \uncertain{facile} (?) \ill{} \ill{} \ill{} \struck{\ill{} \ill{} \ill{} \ill{} $\mathcal{O}_{Y,y_0}$}.
\marginal{\uncertain{cas} \ill{}}

D.
\[
  \dim \mathcal{O}_{X,x} = \underbrace{\dim \mathcal{O}_{Y,y_0}}_{n} + \underbrace{\dim X_1}_{\geq k+1}
\]

\page{26}
Donc divisant $\mathcal{O}_{X,x}$ par un idéal défini par $n$ éléments, on \struck{\ill{}} trouve par le th. général \ill{} (\ill{} : $\mathcal{O}_{X,x}$ \uncertain{anneau} \uncertain{formellement} \ill{} de dim $\geq K+1 = (n+k)+1$] \ill{} \ill{} \ill{} connexe en dim $\geq$ \struck{\ill{}} $K - n = k$, cqfd.
\note{$K$ est une capitale au tracé net ; elle vaut ici $n+k$. La parenthèse ouverte avant « anneau » est fermée par un crochet.}

\emph{Cor. 1} au th.~2. Variante \struck{\ill{}} « géométriquement \ill{} », « \ill{} fibres géométriques.

\emph{Cor. 2} Variante avec $f$ \uncertain{propre} \ill{} ouvert.

\emph{Cor. 3} $X \subset \mathbf{P}^n_k$, \struck{$X$ \ill{}} les composantes de $X$ de dim $\geq k+1$, et $X$ connexe en dim $\geq k$. Alors pour $H$ \uncertain{système} \uncertain{linéaire}, \ill{} \ill{} d'hyperplans $H_1, \ldots, H_\ell$, $X \cap H_1 \cap \cdots \cap H_\ell = Y$ \ill{} condition analogue avec $\ell$ \ill{} $k$.
\marginal{\uncertain{généralisation} au cas de $\alpha$ \ill{} d'un morphisme $X \to \mathbf{P}^r_k$. \uncertain{Pour une} généralisation \ill{} \ill{} \ill{} \ill{} \ill{} \ill{}}
\note{Note marginale gauche, en bas de page, en partie surchargée.}

\section{Couple de Lefschetz-Grauert $(X, Y)$}
\note{Titre écrit par Grothendieck en haut à droite de la couverture grise (page 27) ; « couple de » est écrit au-dessus d'un mot biffé, illisible.}

\page{29}
III \ill{}
\note{En tête de page, « III » souligné, suivi d'un signe surchargé illisible.}

\emph{Introduction}

Dans le \uncertain{§}~\ill{}, nous avons vu que lorsque $X$ est propre sur $Y$ loc.\ noeth., \struck{\ill{} \ill{}} \ill{} \ill{} \ill{}, et si $Y'$ est une partie fermée de $Y$, $X' = f^{-1}(Y')$, la comparaison des foncteurs \struck{$\mathrm{R}^i f$ et $\mathrm{R}^i \hat{f}$ aux} « \ill{} \ill{} » \ill{} \ill{} complètement \ill{} \ill{} de $Y'$, \ill{} $X'$]. Lorsque $Y$ \struck{\ill{}} est le spectre d'un \ill{} \struck{\ill{}} \ill{} \ill{} complet, et $Y' = V(\mathfrak{I})$, \ill{} \ill{}, complétés par un \ill{} \ill{} l'existence (\uncertain{par}~5), s'exprimant \struck{\ill{} \ill{}} \ill{} \ill{} l'équivalence de la catégorie des \struck{\ill{}} Modules cohérents sur $X$ et de la catégorie des Modules cohérents sur $X'_{/Y'}$, \emph{compatible} \ill{} \ill{} \ill{} \ill{} \ill{}.
\note{Les foncteurs lus « $\mathrm{R}^i f$ » et « $\mathrm{R}^i \hat{f}$ » sont soulignés d'un trait qui les biffe peut-être ; « complètement » est d'une lecture incertaine.}

Dans le \ill{} \ill{}, \uncertain{nous} \struck{\ill{}} \ill{} ces résultats \ill{} \ill{} \uncertain{généralisations} \uncertain{à certaines} situations, où $f\colon X' \to X$ \ill{} \ill{} \ill{} [p.~ex.\ $f$ une imm.\ ouverte \ill{}] \ill{} \ill{} \ill{} \ill{} \ill{} \ill{} \struck{\ill{} \ill{}} \ill{}. \struck{\ill{} \ill{} \ill{}} Le \uncertain{résultat} du \ill{} (\uncertain{d'ailleurs} \ill{} \ill{}) \ill{} \ill{} des §~4 et~5, \ill{} \uncertain{utilisant} \struck{des} \ill{} \emph{th.\ de finitude} (N°~1) \ill{} \ill{} (cf N°~2) \struck{\ill{} \ill{} \ill{}} \struck{\ill{}} \ill{} \ill{} \ill{} des \ill{} \struck{s'exprimant} \struck{\ill{} \ill{}} toutes des \ill{} \uncertain{ainsi que}, \uncertain{puis s'exprimant} :
\marginal{des \ill{} \ill{} \ill{} \ill{} \ill{} \ill{} par des \ill{} \ill{} \ill{}, \ill{}}
\note{Bas de page en palimpseste : un long passage encadré est biffé en zigzag, et des renvois (« N° 1 », « N° 2 », « 4 et 5 », ce dernier récrit à l'encre sur un « 4 ») sont intercalés ; on n'en donne que ce qui se lit.}

\page{30}
l'aide de la notion de \emph{\uncertain{profondeur}}.
Les \ill{} résultats obtenus sont l'outil technique principal pour un \ill{} \struck{\ill{}} \ill{} systématique des \struck{\ill{}} propriétés du type de « Lefschetz », \ill{} les propriétés géométriques d'un schéma $X$ (\struck{\ill{}}) à celles de \ill{} \uncertain{sous-schéma} $Y$, \struck{\ill{}} \emph{via} les propriétés du complété formel $X_{/Y}$. \struck{\ill{} \ill{} \ill{} \ill{}} \struck{\ill{} \ill{}} Ils sont \uncertain{applicables} \uncertain{par} exemple (\uncertain{repris} \ill{} \ill{} \ill{} \ill{} \ill{} \ill{}) au cas d'une \uncertain{section} projective \ill{} $X$ \ill{} \ill{} \ill{} \uncertain{hyperplane} \ill{}, moyennant des hypothèses \struck{\ill{} \ill{}} \ill{} sur la \ill{} de $X$. \struck{\ill{} \ill{}} \ill{} \struck{\ill{}} \ill{} \ill{} H.~Grauert, qui \uncertain{avait d'ailleurs} \ill{} (\uncertain{Bourbaki} \ill{}) \struck{\ill{}} 1958, (\uncertain{utilisant} les th.\ de \ill{} \ill{} Serre [\ \ ]]) \struck{\ill{} \ill{} \ill{}}.
\note{Le haut de page est lu d'après les mots soulignés ou encerclés ; les ajouts encerclés « ainsi » et « développement » (lecture incertaine) sont reliés au texte par des traits. Les crochets vides « [ ] » sont de lui : des références à compléter.}

Pour tout $F$ loc.\ libre,
\begin{itemize}
  \item[(i)] $H^i(X, F) \to H^i(\hat{X}, \hat{F})$ \ill{} bijectif [\ill{} $i \leq \operatorname{prof} F - 2$] \struck{\ill{}} \ill{} \struck{\ill{}} \ill{} \ill{} \ill{} \ill{} ;
  \item[(ii)] \ill{} \ill{} \ill{} \ill{} \ill{} $\operatorname{prof}(\mathcal{O}_{\hat{X}}) \geq 2$, \struck{\ill{} \ill{}} \ill{} $X_{/Y}$, Alors \ill{} \ill{} \ill{} « \ill{} » \ill{}, \ill{} \ill{}, \ill{} \ill{} \ill{} [\ill{}] \ill{} \ill{}.
\end{itemize}
\note{La formule de (ii) est surchargée ; le « 2 » et le chapeau sont nets, l'argument de « prof » ne l'est pas.}

\struck{[\ill{} \ill{} \ill{} \ill{} \ill{}]} \struck{\ill{} des Grothendieck} \struck{\ill{} \ill{}} \ill{} \ill{} \ill{}], (\uncertain{p.~ex.}) \ill{} le th.\ de Lefschetz \ill{} \ill{} \ill{} \ill{} \ill{} [\ \ ]).
\marginal{\ill{} Le théorème \ill{} \ill{} \ill{} \ill{} \ill{} \ill{} \ill{} Grauert.}
\note{Note marginale gauche, écrite verticalement, reliée au texte ; on n'en lit que le début et « Grauert ».}

\page{31}
\begin{enumerate}
  \item[1.] Généralisation des th.\ de finitude formelle, \struck{et} des th.\ de comparaison. Application à un th.\ d'existence.
  \item[2.] Le Théorème de finitude locale et ses applications.
  \item[\struck{3.}] \struck{\ill{} de Lefschetz-Grauert. Remarques.}
  \item[\struck{4.}] \struck{\ill{} \ill{} \ill{} \emph{Exemples} \ill{}}
  \item[3.] Cas particulier I : \struck{Étude} \uncertain{Théorie} \ill{} \ill{} locaux.
  \item[4.] Cas particulier 2 : \struck{Étude} \uncertain{Théorie} des \ill{} projectifs.
  \item[5.] Fermés de Lefschetz-Grauert : Remarques, exemples.
  \item[6.] Propriétés liées aux $\pi_0$, $\pi_1$, Pic.
  \item[7.] [Propriétés de permanence]
  \item[8.] Théorèmes de pureté, et théorèmes de Lefschetz pour le $\pi_1$.
  \item[9.] \emph{Anneaux} \struck{\ill{}} \emph{quasi-factoriels} et \emph{th.\ de Lefschetz pour Pic}.
\end{enumerate}
\marginal{utiliser IV, § \ill{} sur les morphismes étales\ldots}
\note{Plan encadré d'un crochet à gauche. Les anciens items 3 et 4 sont enfermés dans un cadre barré de croisillons, avec un « 1 » inscrit par-dessus. « Théorie » remplace « Étude » en interligne aux items 3 et 4 ; un mot encerclé est ajouté au-dessus de « projectifs » (item 4), illisible. Le $\pi_1$ de l'item 6 et le numéro 8 sont encerclés, et la note marginale encadrée « utiliser IV, § … » les relie par deux flèches. Le numéro 9 est écrit sous un numéro encerclé et hachuré.}

\section{Fermés de Lefschetz-Grauert}
\note{Titre souligné en tête de la page 32, de sa main.}

\page{32}
\emph{1) Définitions, remarques}

\emph{Déf.} Soient $X$ un préschéma loc.\ noeth., $Y$ une partie fermée, $n$ un entier. On dit que $Y$ a la propriété $(\mathrm{LG})_n$ dans $X$ \struck{\ill{} \ill{}}, si pour tout $F$ loc.\ libre sur $X$,
\[
  (*) \qquad H^i(X, F) \longrightarrow H^i(X_{/Y}, F_{/Y})
\]
est bijectif pour $i \leq n$, et injectif pour $i = n+1$.
\marginal{(inutile}
\note{« et injectif pour $i = n+1$ » est encerclé, et la marge gauche porte « inutile » précédé d'une parenthèse encerclée : l'ajout est peut-être déclaré inutile (cf.\ la Question de la Remarque~1).}

On dit que $Y$ a la propriété LG \uncertain{faible} si pour tout $F$ loc.\ libre sur $X$, \struck{elle a la propriété $(\mathrm{LG})_0$}
\[
  H^0(X, F) \xrightarrow{\ \sim\ } H^0(X_{/Y}, F_{/Y})
\]
[i.e.\ $F \mapsto F_{/Y}$ est pl.\ fid.\ \ill{} \ill{} Mod.\ loc.\ libres sur $X$ \ill{} ceux sur $X_{/Y}$]. On dit que $Y$ a la propriété LG \struck{\ill{}} \uncertain{stricte} \ill{} si \ill{} voisinage ouvert $U$ de $Y$, \struck{\ill{}} \uncertain{définit} \ill{} \ill{} \ill{} propriété LG, \uncertain{et} \uncertain{si} de plus pour tout $F$ loc.\ libre sur $X_{/Y}$, \ill{} \ill{} isom.\ \ill{} $F_{/Y}$, \ill{} un faisceau \ill{} de la forme \ill{} : $F$ loc.\ libre sur un voisinage $U$ de $Y$.
\note{Le mot écrit au-dessus de « $X$ » à la troisième définition (« faible », lu d'après la Remarque~1 et la page 36) est d'une lecture incertaine ; de même « stricte ».}
\marginal{$\Rightarrow$ Si $\mathfrak{z} \in X$ \struck{\ill{} \ill{}} tel que $\overline{\mathfrak{z}} \cap Y = \emptyset$, alors $\mathfrak{z} \notin \operatorname{Ass} \mathcal{O}_X$}
\marginal{Question : (LG) implique-t-il \struck{\ill{}} $x \in X$ tel que $\overline{x} \cap Y = \emptyset$ $\Rightarrow$ $\operatorname{prof} \mathcal{O}_{X,x} \geq 2$ ???}
\note{Deux notes encadrées en marge gauche ; la lettre de la première est un $\mathfrak{z}$ gothique.}

Donc la propriété \textbf{LG} stricte signifie que les foncteurs $\mathcal{L}(U) \to \mathcal{L}(V)$ [$U \supset V \supset Y$] \uncertain{sont} pleinement fidèles, et
\[
  \varinjlim_{U \supset Y} \mathcal{L}(U) \longrightarrow \mathcal{L}(X_{/Y})
\]
est une équivalence de catégories.
\note{Sous la formule, de sa main : « "Pseudolimite" de catégories », relié par une flèche à $\varinjlim$, et « Catégorie des Modules loc.\ libres sur $U$ », sous $\mathcal{L}(U)$. En marge gauche, une bulle soulignée : « Exemples \ill{} \ill{} \ill{} ».}

\page{33}
\emph{Remarque 1.} La propriété \textbf{LG} \emph{\uncertain{faible}} \emph{\uncertain{équivaut}} à $(\mathrm{LG})_0$. \uncertain{D'autre} \uncertain{part}, il \uncertain{semble} \ill{} que elle implique que
\[
  H^1(X, F) \longrightarrow H^1(X_{/Y}, F_{/Y})
\]
\uncertain{est} injectif, \ill{} \ill{} \uncertain{cas} \ill{} de $\mathcal{O}_X$ \uncertain{par} $F$ \ill{} \ill{} \uncertain{formellement}, \ill{} \ill{}, \ill{} \ill{} est immédiat. [Question : si \ill{} \ill{} \ill{} \ill{} pour $i \leq n$, $\forall\, F$ loc.\ libre, est-il vrai que \ill{} \ill{} \ill{} pour \struck{\ill{}} $i = n+1$, $\forall\, F$ loc.\ libre ?] Oui, prendre une \struck{\ill{}} \ill{} injection $0 \to F \to I \to G \to 0$ et \uncertain{remarquer} \ill{} \ill{} \ill{} \ill{}, \ill{} \ill{} \ill{} \ill{} \ill{} [\ill{} $G$ \ill{}].
\note{La réponse « Oui, … » est écrite en interligne serré ; lecture très incertaine après « injection ».}

\emph{Remarque 2.} Soit $(L_\alpha)$ une \ill{} de Modules loc.\ libres sur $X$, telle que tout Module loc.\ libre sur $X$ \uncertain{soit} \uncertain{quotient} d'une somme directe de Modules $\check{L}_\alpha$, i.e.\ $\simeq$ \uncertain{sous-module} (\uncertain{loc.\ facteur direct}) d'une somme directe de Modules $L_\alpha$ [Ex : $X$ quasi-\uncertain{affine}, $(L_\alpha)$ réduit à $\mathcal{O}_X$ ; $X$ quasi-projectif, $\mathcal{O}(m)$, $m \geq m_0$]. \uncertain{Alors} $(\mathrm{LG})_n$ \uncertain{équivaut} à 2 conditions
\[
  H^i(X, L_\alpha) \simeq H^i(X_{/Y}, L_{\alpha/Y}) \quad \text{pour } i \leq n, \qquad \text{inj pour } i = n+1.
\]
\struck{[N.B. \ill{} \ill{} \ill{} \ill{} \ill{} \ill{} ! ]}. \uncertain{Démontrons} \ill{} récurrence sur $n$, \uncertain{trivial} si $n < -1$, \uncertain{donc} \ill{} \ill{}, supposons \ill{} \ill{}
\marginal{N.B. \ill{} \ill{} \ill{} \ill{} \ill{} \ill{} condition \ill{} injectif.}
\note{Le crochet « [N.B. … ! ] » est biffé par un long zigzag ; la note marginale, oblique, lui est reliée.}

\page{34}
pour les $n' < n$. \struck{Prouvons que} Alors, nous savons déjà que $(*)$ est bijectif pour $i \leq n-1$, injectif pour $i = n$, \ill{} \ill{} \ill{} \ill{} \ill{}. On a une suite exacte
\[
  0 \to F \to L \to G \to 0
\]
\uncertain{où} ($G$ loc.\ libre), $L$ est une somme directe de Modules loc.\ libres \struck{\ill{}} $\simeq L_\alpha$, \ill{} les notations : \ill{} \ill{} \ill{} pour $L$. \uncertain{Considérons} l'homom.\ de \ill{} \ill{}

\begin{tikzcd}[column sep=small, row sep=small, nodes={font=\scriptsize}]
H^i(L) \arrow[r] \arrow[d, "\alpha"'] & H^i(G) \arrow[r] \arrow[d, "\beta"] & H^{i+1}(F) \arrow[r] \arrow[d, "\varphi"] & H^{i+1}(L) \arrow[r] \arrow[d, "\alpha'"] & H^{i+1}(G) \arrow[d, "\beta'"] \\
H^i(\hat{L}) \arrow[r] & H^i(\hat{G}) \arrow[r] & H^{i+1}(\hat{F}) \arrow[r] & H^{i+1}(\hat{L}) \arrow[r] & H^{i+1}(\hat{G})
\end{tikzcd}

\note{Les flèches horizontales sont tracées comme de simples traits ; les flèches verticales $\alpha$ et $\alpha'$ portent un signe $\wr$ à gauche (isomorphismes).}

pour $i = n-1$, \ill{} \ill{} $i+1 = n, n+1$. Dans le premier cas, \ill{} \ill{} \ill{} \ill{} la flèche \uncertain{médiane} est surjective, dans le deuxième qu'elle est injective.

$\alpha$) $i = n-1$, \ill{} \ill{} par hyp.\ de réc.\ que \struck{\ill{} \ill{}} $\beta$ est surjection, \ill{} \ill{} \ill{} \ill{} \ill{}, \uncertain{par} hyp., $\beta'$ \uncertain{est} injection. Par le lemme des~5, $\varphi$ est \struck{\ill{}} surjection (\struck{\ill{}} bijection).
\note{« $\beta'$ est injection » est encerclé et relié à la ligne ; il remplace un passage biffé « \ill{} par l'hyp.\ de réc.\ \ill{} ».}

$\beta$) $i = n$, \ill{} \ill{} \ill{} \struck{\ill{} \ill{}} \struck{\ill{} \ill{} $\beta$ \ill{}} \ill{} que $\alpha$ et $\beta$ sont bijectives, \ill{} $\alpha$ est \uncertain{surj.}, $\beta$ \uncertain{inj.}, de plus $\varphi'$ \ill{} inj.
\note{Dans la dernière ligne, les primes des lettres $\alpha$, $\beta$ ne se distinguent pas ; le dernier symbole est lu $\varphi'$, avec doute.}

\page{35}
\emph{Remarque 3.} La condition LG \ill{} LG \uncertain{effective} \ill{} \uncertain{stable} \ill{} passage de $X$ à un voisinage ouvert de $Y$. Il n'en est pas de même en général des conditions $(\mathrm{LG})_n$, \struck{\ill{} \ill{}} $(\mathrm{LG})_0$ \struck{\ill{} \ill{} \ill{} \ill{} \ill{}} \ill{} \ill{} \ill{} \struck{\ill{} \ill{} \ill{}} $(\mathrm{LG})_0$ \ill{} \uncertain{probablement} \ill{} $=$ \ill{} $(\mathrm{LG})_0$, [\struck{et d'ailleurs}] \ill{} \ill{} \ill{} fait que un Module loc.\ libre $F$ sur un voisinage $U$ de $Y$ n'est \ill{} \ill{} restriction d'un Module loc.\ libre sur \struck{voisinage} $X$. Si cependant tout $F$ \ill{} \ill{} \ill{} un \emph{quotient} d'un Module loc.\ libre sur $X$ [p.~ex.\ si tout Module coh.\ sur $X$ est quotient d'un Module loc.\ libre, p.~ex.\ $X$ quasi-projectif] \emph{alors} $(\mathrm{LG})_0$ pour $(X, Y)$ implique $(\mathrm{LG})_0$ pour tout $(U, Y)$ i.e.\ LG pour $(X, Y)$. \uncertain{En effet} \ill{} \ill{} : \ill{} \ill{} \ill{} $F$ loc.\ libre \emph{sur $X$},
\[
  H^0(U, F) \longrightarrow H^0(U_{/Y}, F_{/Y})
\]
est bijectif, \ill{} \ill{} \ill{} \ill{} \ill{} surjectif, [car $H^0(X, F) \to H^0(U, F) \to H^0(U_{/Y}, F_{/Y})$ \ill{} \ill{}] \uncertain{reste} à vérifier l'\emph{injectivité}, \struck{\ill{}}, \struck{\ill{} \ill{}} \uncertain{qui résulte de}
\note{Deux lignes au début de la remarque sont enfermées dans un cadre barré de deux grandes courbes ; ce qui s'y lit est donné biffé.}

\page{36}
\[
  x \in U, \quad \overline{x}^{(U)} \cap Y = \emptyset \Longrightarrow x \notin \operatorname{Ass} \mathcal{O}_X ,
\]
\ill{} \ill{} \ill{}

\emph{Lemme.} Si $(Y, X)$ satisfait $(\mathrm{LG})_0$, alors pour tout $x \in X$ tel que $Y \cap \overline{x} = \emptyset$, on a $x \notin \operatorname{Ass} \mathcal{O}_X$, \ill{} \ill{} si $X$ est quasi-projectif.

\emph{Dém.} Sinon, il existerait un \uncertain{sous-faisceau} \ill{} $F$ \struck{\ill{}} de $\mathcal{O}_X$ \uncertain{de} \uncertain{support} $\overline{x}$. Donc il existerait une section non nulle de $F(m) \subset \mathcal{O}_X(m)$ [$m$ assez grand], \struck{\ill{}} \ill{} \ill{} \uncertain{induit} une section nulle de $\mathcal{O}_X(m)_{/Y}$, absurde.

\emph{Conclusion pratique} : si $X$ \uncertain{quasi-projectif}, alors
\[
  (\mathrm{LG})_0 \Longleftrightarrow (\mathrm{LG})_{\text{faible}} \Longleftrightarrow (\mathrm{LG})_X \Longrightarrow H^0(X, \mathcal{O}_X(n)) \simeq H^0(X, \mathcal{O}_X(n)_{/Y})
\]
\uncertain{pour} $n \gg 0$.
\note{Conclusion encadrée. L'indice du troisième membre, lu « $X$ », est incertain. Sous le cadre, trois flèches courbes au crayon, chacune marquée « $n$ », renvoient de la dernière implication vers les membres précédents : la réciproque vaut sans doute pour tout $n$.}

\page{37}
\emph{2) Premières propriétés [concernant $\pi_0$, $\pi_1$, Pic]}

\emph{Proposition.} Supposons que
\[
  H^0(X, \mathcal{O}_X) \longrightarrow H^0(X_{/Y}, \mathcal{O}_{X/Y})
\]
soit injectif [bijectif] ; alors $\pi_0(Y) \to \pi_0(X)$ est surjectif [bijectif], \uncertain{en} particulier $(\mathrm{LG})_{-1}$ [\uncertain{resp.}\ $(\mathrm{LG})_0$] implique que $\pi_0(Y) \to \pi_0(X)$ est surjectif [bijectif].

\emph{Proposition.} Supposons que $(X, Y)$ satisfasse \struck{\ill{}} LG (resp.\ LG stricte), \ill{} \ill{} \ill{} \ill{} \ill{} de $Y$. Alors $\pi_1(Y) \to \pi_1(X)$ est surjectif [resp.
\[
  \pi_1(Y) \longrightarrow \varprojlim_{U \text{ voisinage ouvert de } Y} \pi_1(U)
\]
est un homom.\ \struck{surjectif} \struck{[resp.}\ bijectif]].
\note{L'ajout « \ill{} de $Y$ », en haut à droite, est relié à l'énoncé par un trait ; le point-base des $\pi_1$ y est peut-être précisé.}
\marginal{(i) \ill{} \ldots\ (ii) \ill{} \ldots\ Un homom. $F \to G$ \ill{} \ill{} \ill{} \ill{} \ill{} \ill{} \ill{} \ill{} \ill{}}
\note{La marge gauche porte, écrites verticalement et scandées d'accolades, des notes « (i) », « (ii) » sur les homomorphismes de Modules ; presque illisibles.}

Résultats \ill{} \ill{} plus généraux

\struck{Corollaire} \emph{Lemme.} Supposons LG. \struck{Alors \ill{} \ill{}} $F$ loc.\ libre sur $X$, (i) toute structure d'alg.\ \struck{[\ill{} \ill{}, \ill{} \ill{}]} sur $\hat{F}$ provient d'une structure \struck{\ill{}} d'algèbre unique sur $F$, \ill{} \ill{} \ill{} (\ill{} \ill{}, \ill{} \ill{}) \ill{} \ill{} \ill{} ; (ii) les \ill{} de $\hat{F}$ \struck{\ill{}} \ill{} produits de \uncertain{facteurs} \ill{} \ill{} correspondent \ill{} \ill{} \ill{} \ill{} de $\hat{F}$ \ill{} celles de $F$ [\ill{}]. Donc si $X' = \operatorname{Spec}(F)$, \ill{} $X'$ connexe \ill{} \ill{} $\hat{X}'_{/Y}$ connexe
\marginal{[\ill{} \ill{} \ill{} \ill{} \ill{} \ill{} v.l. fidèle \ill{} \ill{} catégories \ill{}]}
\note{Le paragraphe du Lemme est barré de quatre longs traits obliques ; il reste lisible dans ses formules. « Lemme » remplace « Corollaire » biffé.}

\page{38}
\ill{} \ill{} \ill{}, utilisant l'équivalence de la catégorie des rev.\ étales de \struck{$X$} \ill{} et des rev.\ étales de $X_{/Y}$, \ill{}

\emph{Corollaire.} Si (LG) \ill{}, alors le foncteur $X' \mapsto X' \times_X Y$ des rev.\ étales de $X$ dans les rev.\ étales de $Y$ est pleinement fidèle. \struck{Si $X$ (LG) strict \ill{}, \ill{}} Il \ill{} \ill{} \ill{} remplacer $Y$ par un voisinage ouvert \ill{} de $Y$. Si (LG) strict \ill{}, \struck{alors} \ill{} tout rev.\ étale $Y'$ de $Y$, il existe un \struck{\ill{}} voisinage ouvert $U$ de $Y$, et un rev.\ étale $U'$ de $U$, \struck{\ill{}} \ill{}
\[
  U' \times_U Y \simeq Y' .
\]
(\uncertain{Si} $U$ est \uncertain{choisi}) $U'$ est \ill{} \ill{} : \ill{} \ill{} \ill{} \ill{}).

\emph{Remarque.} Nous verrons plus loin des conditions, liées aux « Th.\ de pureté » permettant de conclure que $\pi_1(Y) \to \pi_1(X)$ est \struck{\ill{}} bijectif. Kif-Kif pour Pic \ldots
\note{Dans $\pi_1(Y)$ le $Y$ est récrit sur un $X$ ; « bijectif » est écrit au-dessus d'un mot biffé.}

\page{39}
\emph{Proposition} (i) Si (LG) est satisfaite, alors \ill{} \struck{\ill{} voisinage ouvert} $\operatorname{Pic}(X) \to \operatorname{Pic}(X_{/Y})$ est injectif, \struck{\ill{}} \ill{} plus généralement, \ill{} \ill{} \ill{} \ill{}
\[
  \varinjlim_{U \supset Y} \operatorname{Pic}(U) \longrightarrow \operatorname{Pic}(X_{/Y})
\]
\ill{} \ill{} \ill{} transition injectifs.

(ii) Si (LG) strict est satisfait, \ill{} \ill{} est un isom.
\note{Au (i), le mot biffé devant « voisinage ouvert » est un « tt ».}

\emph{Corollaire.} Supposons que $Y$ « est » un diviseur de Cartier sur $X$, \ill{} $\Lambda \simeq \underline{L}(Y)$ le \uncertain{faisceau} \ill{} \uncertain{associé}, \ill{} $i^*(\Lambda)^{-1} \simeq \mathcal{J}/\mathcal{J}^2$, \uncertain{où} $\mathcal{J}$ est l'Idéal \uncertain{définissant} $Y$. [$i\colon Y \to X$ désigne l'injection canonique.]
\marginal{\ill{} aussi \ill{} \ill{} \ill{} des propriétés}

(i) Supposons $H^1(Y, \Lambda^{-n}) = 0$ pour $n = 1, 2, \ldots$, et (LG). Alors $\operatorname{Pic}(X) \to \operatorname{Pic}(Y)$ est \emph{injectif}. [\ill{} \ill{} les \struck{\ill{} \ill{} \ill{}} $\operatorname{Pic}(U) \to \operatorname{Pic}(Y)$ \ill{} injectifs, \struck{\ill{} \ill{} \ill{}} \ill{}]

(ii) Supposons \struck{$H^1(Y, \Lambda^{-n}) = 0$} et $H^2(Y, \Lambda^{-n}) = 0$ pour $n = 1, 2, \ldots$, et (LG) strict. Alors
\[
  \varinjlim_{U \ni Y} \operatorname{Pic}(U) \xrightarrow{\ \sim\ } \operatorname{Pic}(Y) .
\]
\marginal{[\ill{} \ill{} \ill{}, $\varinjlim_{U \supset Y} \operatorname{Pic}(U) \to \operatorname{Pic}(Y)$ \ill{} \ill{} \ill{} injectifs]}
\note{Au (ii), un trait traverse « $H^1(Y, \Lambda^{-n}) = 0$ » ; il n'est pas sûr que ce soit une biffure, la conclusion demandant les deux annulations (cf.\ le Lemme de la page 40). La note marginale gauche, reliée par un trait à (i), rappelle l'injectivité des $\operatorname{Pic}(U) \to \operatorname{Pic}(Y)$.}

\page{40}
\ill{} \ill{} \ill{}
\struck{\ill{} \ill{} \ill{}} \ill{} \ill{} \struck{\ill{}} \ill{} Corollaire de la \ill{} \ill{} \ill{}
\note{Deux lignes de liaison en tête de page, récrites en interligne ; on y lit seulement « Corollaire ».}

\emph{Lemme} (i) Si $H^1(Y, \Lambda^{-n}) = 0$ pour $n = 1, 2, \ldots$, alors \struck{\ill{} Pic} le foncteur $\mathcal{F} \mapsto \mathcal{F}_0$ des \ill{} \struck{\ill{}} inversibles sur $\hat{X}$ \ill{} les modules inversibles sur $Y$ est pl.\ fidèle, \ill{} $\operatorname{Pic}(\hat{X}) \to \operatorname{Pic}(Y)$ est injectif.

(ii) Si de plus, $H^2(Y, \Lambda^{-n}) = 0$ pour $n = 1, 2, \ldots$, alors le foncteur précédent est une équiv.\ de catégories, donc $\operatorname{Pic}(\hat{X}) \to \operatorname{Pic}(Y)$ est un isom.

\emph{3) Théorème de permanence.}
\note{Tout ce qui suit sur la page est annulé par de longs traits obliques croisés ; on le donne tel qu'il se lit.}

\emph{Théorème 1.} Supposons que $(X, Y)$ \struck{$X, Y$} satisfasse $(\mathrm{LG})_n$ et $X$ \ill{} \uncertain{projectif}. Alors \ill{} \ill{} \struck{$(\mathrm{LG})_n$ et} \struck{soit} $f\colon X' \to X$ \ill{} morphisme projectif et plat, \struck{Alors} $f\colon X' \to X$, \ill{}
\[
  Y' = f^{-1}(Y) ,
\]
$(X', Y')$ satisfait $(\mathrm{LG})_n$.

\struck{\emph{Théorème} \ill{} Si $(X, Y)$ satisfait (LG), alors}

\emph{Th.~2} Si $(X, Y)$ satisfait (LG) [resp.\ (LG) \ill{}], \ill{} $X$ \uncertain{quasi-projectif}, --- alors
\note{L'énoncé s'interrompt au bas de la page 40 ; la suite est dans le lot suivant.}

\end{document}
